% !TEX spellcheck = en_US



% ~~~~~~~~~~~~~~~~~~~~~~~~~~~~~~~~~~~~~~~ V
% (NC) 2026 Haluk Bingol
% github.com/halukbingol/LaTeX-Templates
%
% Licensees may copy, distribute, display, and perform the work and make derivative works 
% and remixes based on it only for non-commercial purposes. 
% ~~~~~~~~~~~~~~~~~~~~~~~~~~~~~~~~~~~~~~~ A




% =====================================================================
%  Queues -- Sophomore Level: Interfaces, Generics and Tests
%  ---------------------------------------------------------------------
%  Built on hblib/Template/hblibTemplate.tex with the libraries in ../../hblib/.
%  Programs and their real outputs are in codeoutput/:
%     codeoutput/<Name>.java  a program (classes without main() are only compiled)
%     codeoutput/<Name>.txt   what it printed, made by:  bash run_examples.sh
% =====================================================================

\documentclass[aspectratio=169,10pt,compress]{beamer}

% ~~~~~~~~~~~~~~~~~~~~~~~~~~~~~~~~~~~~~~~~~~~~~~~~~~~~~~~~~~~~~~~~~ V hblib
	\usepackage{../../hblib/hblibPresentation} % lib presentation: theme, i/N, \hSection, algorithm2e, ...
	\usepackage{../../hblib/hblibCodeoutput} % lib code + output: \lstcode, \lstterminal, ...
	\usepackage{../../hblib/hblibDatastructures} % lib data structures: \hString, \hStack, \hQueue
% ~~~~~~~~~~~~~~~~~~~~~~~~~~~~~~~~~~~~~~~~~~~~~~~~~~~~~~~~~~~~~~~~~ A hblib


% xcolor, array (and the L{width} column), listings, tikz and algorithm2e come from the hblib libraries
\usepackage[T1]{fontenc}
\usepackage{lmodern}
\usepackage{booktabs}

\newcommand{\win}{\textbf{Windows}}
\newcommand{\mac}{\textbf{macOS}}

% listings: fixed character cells, so the columns of a program's output stay aligned
\lstset{columns=fixed}

% algorithm2e set-up (the package is loaded by hblibPresentation)
\SetAlCapFnt{\small}
\SetAlCapNameFnt{\small}
\SetKwProg{Fn}{function}{}{end}
\SetKw{KwError}{error}

% ---------------------------------------------------------------------
\title{Queues}
\subtitle{Sophomore Level: Interfaces, Generics and Tests}
\author[Bingol]{Haluk O. Bingol}

\institute[FCIS]{
    Faculty of Computer and Information Sciences\\
    Yeditepe University
}
\date{
	{\tiny \hbTimeStamp}
}

\begin{document}




% ~~~~~~~~~~~~~~~~~~~~~~~~~~~~~~~~~~~~~~ V frm
\begin{frame}[t,fragile]
  \titlepage
\end{frame}
% ~~~~~~~~~~~~~~~~~~~~~~~~~~~~~~~~~~~~~~ A frm




% ~~~~~~~~~~~~~~~~~~~~~~~~~~~~~~~~~~~~~~ V frm
\begin{frame}[t,fragile] \frametitle{Outline}
  \begin{columns}[T]
    \begin{column}{0.48\textwidth}
      \tableofcontents[sections={1-3}, hideallsubsections]
    \end{column}
    \begin{column}{0.48\textwidth}
      \tableofcontents[sections={4-5}, hideallsubsections]
    \end{column}
  \end{columns}
\end{frame}
% ~~~~~~~~~~~~~~~~~~~~~~~~~~~~~~~~~~~~~~ A frm

% ~~~~~~~~~~~~~~~~~~~~~~~~~~~~~~~~~~~~~~ V
\hSection[Design]{Interface and Implementation}
% ~~~~~~~~~~~~~~~~~~~~~~~~~~~~~~~~~~~~~~ A




% ~~~~~~~~~~~~~~~~~~~~~~~~~~~~~~~~~~~~~~ V frm
\begin{frame}[t,fragile] \frametitle{Separate \emph{what} from \emph{how}}
  \begin{columns}[T]
    \begin{column}{0.46\textwidth}
      \lstcode[basicstyle=\ttfamily\tiny]{codeoutput/QueueADT.java}{Java}
      \vspace{0.2em}
      \scriptsize \texttt{QueueUnderflowException} is a five-line class that extends \texttt{RuntimeException}.
    \end{column}
    \begin{column}{0.5\textwidth}
      \begin{itemize}\small
        \item The \textbf{interface} is the ADT: the operations, nothing about arrays.
        \item Callers write \texttt{QueueADT<String> q = new ArrayQueue<>()}; switching to a linked implementation changes \textbf{one} line.
        \item \texttt{extends Iterable<T>}: every queue works in a for-each loop.
        \item An \textbf{unchecked} exception (\texttt{RuntimeException}): dequeuing an empty queue is a programming error, like an array index error.
      \end{itemize}
    \end{column}
  \end{columns}
\end{frame}
% ~~~~~~~~~~~~~~~~~~~~~~~~~~~~~~~~~~~~~~ A frm




% ~~~~~~~~~~~~~~~~~~~~~~~~~~~~~~~~~~~~~~ V frm
\begin{frame}[t,fragile] \frametitle{\texttt{ArrayQueue<T>}: the implementation}
  \begin{columns}[T]
    \begin{column}{0.49\textwidth}
      \lstcode[basicstyle=\ttfamily\tiny, firstline=1, lastline=21, firstnumber=1]{codeoutput/ArrayQueue.java}{Java}
    \end{column}
    \begin{column}{0.49\textwidth}
      \lstcode[basicstyle=\ttfamily\tiny, firstline=22, lastline=41, firstnumber=22]{codeoutput/ArrayQueue.java}{Java}
    \end{column}
  \end{columns}
  \small The class says \texttt{implements QueueADT<T>}: the compiler checks that every method of the interface exists.
\end{frame}
% ~~~~~~~~~~~~~~~~~~~~~~~~~~~~~~~~~~~~~~ A frm




% ~~~~~~~~~~~~~~~~~~~~~~~~~~~~~~~~~~~~~~ V frm
\begin{frame}[t,fragile] \frametitle{Program to the interface}
  \begin{columns}[T]
    \begin{column}{0.46\textwidth}
      \begin{itemize}\small
        \item The variable type is \texttt{QueueADT<String>}, not \texttt{ArrayQueue<String>}.
        \item The rest of the program does not know how the queue is stored.
        \item A linked implementation, a test double, or \texttt{ArrayDeque} can be substituted later.
      \end{itemize}
      \vspace{0.3em}


% +++++++++++++++++++++++++++++++++++++++ V lst
%: -Lst.
\begin{lstlisting}[style=codesnippet, language=Java, basicstyle=\ttfamily\scriptsize]
QueueADT<String> q = new ArrayQueue<>();
q.enqueue("Ali");
System.out.println(q.dequeue());
\end{lstlisting}
% +++++++++++++++++++++++++++++++++++++++ A lst


    \end{column}
    \begin{column}{0.5\textwidth}
      \lstoutput[basicstyle=\ttfamily\tiny]{codeoutput/UseQueue.txt}
    \end{column}
  \end{columns}
\end{frame}
% ~~~~~~~~~~~~~~~~~~~~~~~~~~~~~~~~~~~~~~ A frm

% ~~~~~~~~~~~~~~~~~~~~~~~~~~~~~~~~~~~~~~ V
\hSection[Generics]{Generics and Arrays}
% ~~~~~~~~~~~~~~~~~~~~~~~~~~~~~~~~~~~~~~ A




% ~~~~~~~~~~~~~~~~~~~~~~~~~~~~~~~~~~~~~~ V frm
\begin{frame}[t,fragile] \frametitle{Why not \texttt{new T[n]}?}
  \begin{columns}[T]
    \begin{column}{0.5\textwidth}
      \small Java does not know the type \texttt{T} when it runs: it is \textbf{erased}. So it cannot create an array of \texttt{T}.


% +++++++++++++++++++++++++++++++++++++++ V lst
%: -Lst.
\begin{lstlisting}[style=codesnippet, language=Java, basicstyle=\ttfamily\scriptsize]
T[] a = new T[n];              // error
T[] a = (T[]) new Object[n];   // warning
\end{lstlisting}
% +++++++++++++++++++++++++++++++++++++++ A lst


      \vspace{0.3em}
      \small The cast is \textbf{unchecked}: the compiler cannot prove it. We suppress the warning in one place.
    \end{column}
    \begin{column}{0.46\textwidth}
      \begin{block}{A heap of Object}
        \small The array really is an \texttt{Object[]}. Every element put in is a \texttt{T}, so every element taken out is a \texttt{T} too.
      \end{block}
      \begin{block}{Danger}
        \small If the array escapes as a \texttt{T[]}, code could store a wrong type in it. Keep it \texttt{private} and never return it.
      \end{block}
    \end{column}
  \end{columns}
\end{frame}
% ~~~~~~~~~~~~~~~~~~~~~~~~~~~~~~~~~~~~~~ A frm




% ~~~~~~~~~~~~~~~~~~~~~~~~~~~~~~~~~~~~~~ V frm
\begin{frame}[t,fragile] \frametitle{The cast is safe here, but only inside}
  \begin{columns}[T]
    \begin{column}{0.52\textwidth}
      \lstcode[basicstyle=\ttfamily\tiny, firstline=43, lastline=51, firstnumber=43]{codeoutput/ArrayQueue.java}{Java}
    \end{column}
    \begin{column}{0.44\textwidth}
      \begin{itemize}\small
        \item \texttt{resize} creates one \texttt{Object[]} and copies into it.
        \item A private method cannot be called from outside with a wrong type.
        \item \texttt{@SuppressWarnings("unchecked")} is on the smallest possible scope: one line, one method.
      \end{itemize}
    \end{column}
  \end{columns}
\end{frame}
% ~~~~~~~~~~~~~~~~~~~~~~~~~~~~~~~~~~~~~~ A frm




% ~~~~~~~~~~~~~~~~~~~~~~~~~~~~~~~~~~~~~~ V frm
\begin{frame}[t,fragile] \frametitle{Iterating from the front}
  \begin{columns}[T]
    \begin{column}{0.52\textwidth}
      \lstcode[basicstyle=\ttfamily\tiny, firstline=53, lastline=63, firstnumber=53]{codeoutput/ArrayQueue.java}{Java}
    \end{column}
    \begin{column}{0.44\textwidth}
      \small The iterator keeps its own counter and follows \texttt{head} around the array, so the order is front to rear even when the queue has wrapped.\\[0.3em]
      \uncover<2->{\lstoutput[basicstyle=\ttfamily\tiny]{codeoutput/IterateDemo.txt}}
    \end{column}
  \end{columns}
\end{frame}
% ~~~~~~~~~~~~~~~~~~~~~~~~~~~~~~~~~~~~~~ A frm

% ~~~~~~~~~~~~~~~~~~~~~~~~~~~~~~~~~~~~~~ V
\hSection[Memory]{Resizing and Memory}
% ~~~~~~~~~~~~~~~~~~~~~~~~~~~~~~~~~~~~~~ A

% ~~~~~~~~~~~~~~~~~~~~~~~~~~~~~~~~~~~~~~ V
\subsection[Copies]{How much copying?}
% ~~~~~~~~~~~~~~~~~~~~~~~~~~~~~~~~~~~~~~ A




% ~~~~~~~~~~~~~~~~~~~~~~~~~~~~~~~~~~~~~~ V frm
\begin{frame}[t,fragile] \frametitle{Growing and shrinking}
  \lstcoderightstep[0.55][basicstyle=\ttfamily\tiny, firstline=43, lastline=51, firstnumber=43][basicstyle=\ttfamily\tiny]{codeoutput/ArrayQueue.java}{Java}{codeoutput/GrowDemo.txt}
  \small One \texttt{resize} method serves both grow and shrink: copy in queue order, then \texttt{head = 0}, \texttt{tail = size}.
\end{frame}
% ~~~~~~~~~~~~~~~~~~~~~~~~~~~~~~~~~~~~~~ A frm




% ~~~~~~~~~~~~~~~~~~~~~~~~~~~~~~~~~~~~~~ V frm
\begin{frame}[t,fragile] \frametitle{Amortized cost of doubling}
  Enqueue $n$ elements into a queue that starts with capacity 1 and doubles:
  \[
  	\text{copies} = 1 + 2 + 4 + \dots + 2^{k} = 2^{k+1} - 1 < 2n
  \]
  \[
  	\text{total work} < n + 2n = 3n
  \]
  \begin{itemize}\small
    \item The average cost of one enqueue is less than 3 steps: $O(1)$ \hRemark{amortized}.
    \item The single enqueue that copies still costs $O(n)$.
    \item Shrinking at \textbf{one quarter} (not one half) means an enqueue/dequeue pair at the boundary does not copy every time.
  \end{itemize}
\end{frame}
% ~~~~~~~~~~~~~~~~~~~~~~~~~~~~~~~~~~~~~~ A frm

% ~~~~~~~~~~~~~~~~~~~~~~~~~~~~~~~~~~~~~~ V
\subsection[Loitering]{Loitering}
% ~~~~~~~~~~~~~~~~~~~~~~~~~~~~~~~~~~~~~~ A




% ~~~~~~~~~~~~~~~~~~~~~~~~~~~~~~~~~~~~~~ V frm
\begin{frame}[t,fragile] \frametitle{Loitering: dequeued, but not garbage}
  \begin{columns}[T]
    \begin{column}{0.5\textwidth}
      \begin{block}{With \texttt{a[head] = null}}
        \lstoutput[basicstyle=\ttfamily\tiny]{codeoutput/Loitering.txt}
      \end{block}
    \end{column}
    \begin{column}{0.46\textwidth}
      \begin{block}{Without it}
        \lstoutput[basicstyle=\ttfamily\tiny]{codeoutput/Leak.txt}
      \end{block}
    \end{column}
  \end{columns}
  \vspace{0.4em}
  \begin{itemize}\small
    \item The array still holds a reference to a dequeued object.
    \item The garbage collector cannot free it, although the queue no longer uses it: \hRemark{loitering}.
    \item The fix costs one assignment. It matters for queues of large objects and for long-running programs.
  \end{itemize}
\end{frame}
% ~~~~~~~~~~~~~~~~~~~~~~~~~~~~~~~~~~~~~~ A frm

% ~~~~~~~~~~~~~~~~~~~~~~~~~~~~~~~~~~~~~~ V
\hSection[Testing]{Testing a Queue}
% ~~~~~~~~~~~~~~~~~~~~~~~~~~~~~~~~~~~~~~ A




% ~~~~~~~~~~~~~~~~~~~~~~~~~~~~~~~~~~~~~~ V frm
\begin{frame}[t,fragile] \frametitle{Tests without a framework}
  \begin{columns}[T]
    \begin{column}{0.52\textwidth}
      \lstcode[basicstyle=\ttfamily\tiny, lastline=20]{codeoutput/TestArrayQueue.java}{Java}
    \end{column}
    \begin{column}{0.44\textwidth}
      \lstcode[basicstyle=\ttfamily\tiny, firstline=21, firstnumber=21]{codeoutput/TestArrayQueue.java}{Java}
    \end{column}
  \end{columns}
  \small One line per check, and a final count. A test that fails tells you which rule broke, not only that something broke.
\end{frame}
% ~~~~~~~~~~~~~~~~~~~~~~~~~~~~~~~~~~~~~~ A frm




% ~~~~~~~~~~~~~~~~~~~~~~~~~~~~~~~~~~~~~~ V frm
\begin{frame}[t,fragile] \frametitle{The tests pass}
  \lstoutput[basicstyle=\ttfamily\scriptsize]{codeoutput/TestArrayQueue.txt}
  \vspace{0.3em}
  \small The most valuable cases are the edges: a new queue, a full queue, a wrapped queue, and a dequeue on an empty queue.
\end{frame}
% ~~~~~~~~~~~~~~~~~~~~~~~~~~~~~~~~~~~~~~ A frm




% ~~~~~~~~~~~~~~~~~~~~~~~~~~~~~~~~~~~~~~ V frm
\begin{frame}[t,fragile] \frametitle{The same tests with JUnit 5}
  \begin{columns}[T]
    \begin{column}{0.56\textwidth}


% +++++++++++++++++++++++++++++++++++++++ V lst
%: -Lst.
\begin{lstlisting}[style=codesnippet, language=Java, basicstyle=\ttfamily\scriptsize]
import org.junit.jupiter.api.Test;
import static org.junit.jupiter.api.Assertions.*;

class ArrayQueueTest {
    @Test void fifoOrder() {
        ArrayQueue<Integer> q = new ArrayQueue<>();
        q.enqueue(1); q.enqueue(2);
        assertEquals(1, q.dequeue());
        assertEquals(2, q.dequeue());
    }
    @Test void underflow() {
        assertThrows(QueueUnderflowException.class,
                     () -> new ArrayQueue<>().dequeue());
    }
}
\end{lstlisting}
% +++++++++++++++++++++++++++++++++++++++ A lst


    \end{column}
    \begin{column}{0.4\textwidth}
      \begin{itemize}\small
        \item \texttt{assertEquals} compares the value and prints both when they differ.
        \item \texttt{assertThrows} checks that the right exception is thrown.
        \item One test method per rule; the name says what is checked.
      \end{itemize}
      \vspace{0.3em}
      \hRemark{Not run here}: it needs \texttt{junit-jupiter} on the class path (Maven or Gradle adds it).
    \end{column}
  \end{columns}
\end{frame}
% ~~~~~~~~~~~~~~~~~~~~~~~~~~~~~~~~~~~~~~ A frm

% ~~~~~~~~~~~~~~~~~~~~~~~~~~~~~~~~~~~~~~ V
\hSection[Library]{The Library: \texttt{Queue} and \texttt{Deque}}
% ~~~~~~~~~~~~~~~~~~~~~~~~~~~~~~~~~~~~~~ A




% ~~~~~~~~~~~~~~~~~~~~~~~~~~~~~~~~~~~~~~ V frm
\begin{frame}[t,fragile] \frametitle{The methods}
  \begin{center}\footnotesize
  \begin{tabular}{@{}lll@{}}
    \toprule
    \textbf{Task} & \textbf{throws on failure} & \textbf{returns a value on failure} \\
    \midrule
    add at the rear   & \texttt{add(e)}      & \texttt{offer(e)} \\
    remove the front  & \texttt{remove()}    & \texttt{poll()} \\
    look at the front & \texttt{element()}   & \texttt{peek()} \\
    \midrule
    add at the front  & \texttt{addFirst(e)} & \texttt{offerFirst(e)} \\
    remove the rear   & \texttt{removeLast()} & \texttt{pollLast()} \\
    look at the rear  & \texttt{getLast()}   & \texttt{peekLast()} \\
    \bottomrule
  \end{tabular}
  \end{center}
  \vspace{0.4em}
  \begin{itemize}\small
    \item \texttt{Queue} has the first three rows; \texttt{Deque} has all of them.
    \item \texttt{ArrayDeque} implements \texttt{Deque}: it can be a queue, a stack, or both.
    \item \texttt{LinkedList} also implements \texttt{Deque}, but each element costs a node.
  \end{itemize}
\end{frame}
% ~~~~~~~~~~~~~~~~~~~~~~~~~~~~~~~~~~~~~~ A frm




% ~~~~~~~~~~~~~~~~~~~~~~~~~~~~~~~~~~~~~~ V frm
\begin{frame}[t,fragile] \frametitle{Both ends at once}
  \lstcoderight[0.52][basicstyle=\ttfamily\tiny][basicstyle=\ttfamily\scriptsize]{codeoutput/DequeDemo.java}{Java}{codeoutput/DequeDemo.txt}
  \small \texttt{addFirst} + \texttt{pollFirst} is a stack; \texttt{addLast} + \texttt{pollFirst} is a queue.
\end{frame}
% ~~~~~~~~~~~~~~~~~~~~~~~~~~~~~~~~~~~~~~ A frm

% ~~~~~~~~~~~~~~~~~~~~~~~~~~~~~~~~~~~~~~ V
\hSection[Summary]{Summary}
% ~~~~~~~~~~~~~~~~~~~~~~~~~~~~~~~~~~~~~~ A




% ~~~~~~~~~~~~~~~~~~~~~~~~~~~~~~~~~~~~~~ V frm
\begin{frame}[t,fragile] \frametitle{Recipes}
  \begin{center}\small
  \begin{tabular}{@{}L{0.28\textwidth}L{0.62\textwidth}@{}}
    \toprule
    \textbf{Need} & \textbf{Do} \\
    \midrule
    an ADT for callers      & a generic \texttt{interface QueueADT<T>} \\
    a queue in an array     & two indices, \texttt{\% n}, and a \texttt{size} counter \\
    to grow                 & double, copy in queue order, \texttt{head = 0}, \texttt{tail = size} \\
    to avoid loitering      & \texttt{a[head] = null} before moving \texttt{head} \\
    to test                 & empty, full, wrapped, underflow, 100-element order \\
    a queue from the library& \texttt{Deque<T> q = new ArrayDeque<>()} \\
    \bottomrule
  \end{tabular}
  \end{center}
\end{frame}
% ~~~~~~~~~~~~~~~~~~~~~~~~~~~~~~~~~~~~~~ A frm




% ~~~~~~~~~~~~~~~~~~~~~~~~~~~~~~~~~~~~~~ V frm
\begin{frame}[t,fragile] \frametitle{Exercises}
  \begin{enumerate}\small
    \item Add \texttt{clear()} and \texttt{contains(T x)} to \texttt{ArrayQueue}.
    \item Write \texttt{LinkedQueue<T>} that implements the same \texttt{QueueADT<T>}. Change one line in \texttt{UseQueue}.
    \item Make \texttt{ArrayQueue} shrink at one quarter and test the boundary: enqueue, dequeue, enqueue, dequeue, \dots
    \item Write a JUnit test for a wrapped iterator (enqueue 5, dequeue 2, enqueue 3, then check the order).
    \item Measure how many copies 1000 enqueues cause with doubling and with growth by 1.
  \end{enumerate}
\end{frame}
% ~~~~~~~~~~~~~~~~~~~~~~~~~~~~~~~~~~~~~~ A frm




% ~~~~~~~~~~~~~~~~~~~~~~~~~~~~~~~~~~~~~~ V frm
\begin{frame}[t,fragile]
	\begin{center}
		\vspace{4cm}
		\hRemark{\Huge Questions?}
	\end{center}
\end{frame}
% ~~~~~~~~~~~~~~~~~~~~~~~~~~~~~~~~~~~~~~ A frm




\end{document}
